\section*{ROOT SYSTEMS}

\begin{enumerate}
\def\labelenumi{\arabic{enumi})}
\setcounter{enumi}{4}
\item
  For any element \(\alpha\) of the dual \(\mathfrak{h}^*\) of \(\mathfrak{h}\) , let \(\mathfrak{g}^{\alpha}\) be the set of \(x\in \mathfrak{g}\) such that \([h,x] = \alpha (h)x\) for all \(h\in \mathfrak{h}\) . If \(\alpha = 0\) , then \(\mathfrak{g}^{\alpha} = \mathfrak{h}\) . Any \(\alpha \in \mathfrak{h}^{*} - \{0\}\) such that \(\mathfrak{g}^{\alpha}\neq 0\) is called a root of \((\mathfrak{g},\mathfrak{h})\) . Denote by \(\mathrm{R}(\mathfrak{g},\mathfrak{h})\) (or simply by R) the set of roots of \((\mathfrak{g},\mathfrak{h})\) . This is a reduced root system in \(\mathfrak{h}^*\) in the sense of Chap. VI, §1, no. 4. The algebra \(\mathfrak{g}\) is simple if and only if R is irreducible.
\item
  For all \(\alpha \in \mathbb{R}\) , \(\mathfrak{g}^{\alpha}\) is of dimension 1. The vector space \([\mathfrak{g}^{\alpha}, \mathfrak{g}^{-\alpha}]\) is contained in \(\mathfrak{h}\) , is of dimension 1, and contains a unique element \(H_{\alpha}\) such that \(\alpha(H_{\alpha}) = 2\) ; we have \(H_{\alpha} = \alpha^{\vee}\) (Chap. VI, §1, no. 1); the set of \(H_{\alpha}\) , for \(\alpha \in \mathbb{R}\) , is the inverse root system \(\mathbb{R}^{\vee}\) of \(\mathbb{R}\) .
\item
  We have \(\mathfrak{g} = \mathfrak{h} \oplus \bigoplus_{\alpha \in \mathrm{R}} \mathfrak{g}^{\alpha}\) . There exists a family \((X_{\alpha})_{\alpha \in \mathrm{R}}\) such that, for all \(\alpha \in \mathrm{R}\) , we have \(X_{\alpha} \in \mathfrak{g}^{\alpha}\) and \([X_{\alpha}, X_{-\alpha}] = -H_{\alpha}\) . Every \(x \in \mathfrak{g}\) can be written uniquely in the form
\end{enumerate}

\[
x = h + \sum_ {\alpha \in \mathrm{R}} \lambda_ {\alpha} X _ {\alpha}, \quad \text { where } h \in \mathfrak {h}, \lambda_ {\alpha} \in k.
\]

The bracket of two elements can be calculated by means of the formulas

\[
\begin{array}{c} [ h, X _ {\alpha} ] = \alpha (h) X _ {\alpha} \\ \relax [ X _ {\alpha}, X _ {\beta} ] = 0 \quad \text {if} \alpha + \beta \notin \mathrm{R} \cup \{0 \} \\ \relax [ X _ {\alpha}, X _ {- \alpha} ] = - H _ {\alpha} \\ \relax [ X _ {\alpha}, X _ {\beta} ] = \mathrm{N} _ {\alpha \beta} X _ {\alpha + \beta} \quad \text {if} \alpha + \beta \in \mathrm{R}, \end{array}
\]

the \(\mathrm{N}_{\alpha \beta}\) being non-zero elements of \(k\) .

\begin{enumerate}
\def\labelenumi{\arabic{enumi})}
\setcounter{enumi}{7}
\tightlist
\item
  Let B be a basis of R. The algebra g is generated by the \(X_{\alpha}\) and the \(X_{-\alpha}\) for \(\alpha \in B\) . We have \([X_{\alpha}, X_{-\beta}] = 0\) if \(\alpha, \beta \in B\) and \(\alpha \neq \beta\) . Let \((n(\alpha, \beta))_{\alpha, \beta \in B}\) be the Cartan matrix of R (relative to B). We have \(n(\alpha, \beta) = \alpha(H_{\beta})\) . If \(\alpha, \beta \in B\) and \(\alpha \neq \beta\) , \(n(\alpha, \beta)\) is a negative integer and
\end{enumerate}

\[
(\operatorname{ad} X _ {\beta}) ^ {1 - n (\alpha , \beta)} X _ {\alpha} = 0 \quad \text { and } \quad (\operatorname{ad} X _ {- \beta}) ^ {1 - n (\alpha , \beta)} X _ {- \alpha} = 0.
\]

\begin{enumerate}
\def\labelenumi{\arabic{enumi})}
\setcounter{enumi}{8}
\tightlist
\item
  If \(\alpha, \beta, \alpha + \beta \in \mathbb{R}\) , let \(q_{\alpha \beta}\) be the largest integer \(j\) such that \(\beta - j\alpha \in \mathbb{R}\) . The family \((X_{\alpha})_{\alpha \in \mathbb{R}}\) in 7) can be chosen so that \(\mathrm{N}_{\alpha, \beta} = \mathrm{N}_{-\alpha, -\beta}\) if \(\alpha, \beta, \alpha + \beta \in \mathbb{R}\) . Then \(\mathrm{N}_{\alpha \beta} = \pm (q_{\alpha \beta} + 1)\) . There exists an involutive automorphism \(\theta\) of \(\mathfrak{g}\) that takes \(X_{\alpha}\) to \(X_{-\alpha}\) for all \(\alpha \in \mathbb{R}\) ; we have \(\theta(h) = -h\) for all \(h \in \mathfrak{h}\) . The \(\mathbf{Z}\) -submodule \(\mathfrak{g}_{\mathbf{Z}}\) of \(\mathfrak{g}\) generated by the \(H_{\alpha}\) and the \(X_{\alpha}\) is a \(\mathbf{Z}\) -Lie subalgebra of \(\mathfrak{g}\) , and the canonical map \(\mathfrak{g}_{\mathbf{Z}} \otimes_{\mathbf{Z}} k \to \mathfrak{g}\) is an isomorphism.
\end{enumerate}

The pair \((\mathfrak{g},\mathfrak{h})\) can be obtained by extension of scalars from a split semi-simple Q-Lie algebra.

\begin{enumerate}
\def\labelenumi{\arabic{enumi})}
\setcounter{enumi}{9}
\item
  The Weyl group, group of weights, \ldots{} of R is called the Weyl group, group of weights, \ldots{} of \((\mathfrak{g},\mathfrak{h})\) . The Weyl group will be denoted by W in what follows. We consider its operation, not only on \(h^{*}\) , but also on h (by transport of structure). If \(h_{Q}\) (resp. \(h_{Q}^{*}\) ) denotes the Q-vector subspace of h (resp. \(h^{*}\) ) generated by the \(H_{\alpha}\) (resp. the \(\alpha\) ), then h (resp. \(h^{*}\) ) can be canonically identified with \(h_{Q} \otimes_{Q} k\) (resp. \(h_{Q}^{*} \otimes_{Q} k\) ), and \(h_{Q}^{*}\) can be identified with the dual of \(h_{Q}\) . When we speak of the Weyl chambers of R, these are understood to be in \(h_{R} = h_{Q} \otimes_{Q} R\) or \(h_{R}^{*} = h_{Q}^{*} \otimes_{Q} R\) .
\item
  Let \(\varPhi\) be the Killing form of \(\mathfrak{g}\) . If \(\alpha + \beta \neq 0\) , \(\mathfrak{g}^{\alpha}\) and \(\mathfrak{g}^{\beta}\) are orthogonal with respect to \(\varPhi\) . The restriction of \(\varPhi\) to \(\mathfrak{g}^{\alpha} \times \mathfrak{g}^{-\alpha}\) is non-degenerate. If \(x, y \in \mathfrak{h}\) , then \(\varPhi(x, y) = \sum_{\alpha \in \mathrm{R}} \alpha(x)\alpha(y)\) . We have \(\varPhi(H_{\alpha}, H_{\beta}) \in \mathbf{Z}\) . The restriction of \(\varPhi\) to \(\mathfrak{h}\) is non-degenerate and invariant under W; its restriction to \(\mathfrak{h}_{\mathbf{Q}}\) is positive.
\item
  The root system of \((\mathfrak{g},\mathfrak{h})\) depends, up to isomorphism, only on g and not on h. By abuse of language, the Weyl group, group of weights, \ldots{} of \((\mathfrak{g},\mathfrak{h})\) are also called the Weyl group, group of weights, \ldots{} of g.
\end{enumerate}

If \(R_{1}\) is a reduced root system, there exists a split semi-simple Lie algebra \((\mathfrak{g}_{1},\mathfrak{h}_{1})\) such that \(\mathrm{R}(\mathfrak{g}_{1},\mathfrak{h}_{1})\) is isomorphic to \(R_{1}\) ; it is unique, up to isomorphism.

The classification of splittable semi-simple Lie algebras is thus reduced to that of root systems.

